在當今城市交通智慧化與綠色出行理念深入發展的背景下，共用單車系統已成為解決“最後一公里”難題、緩解交通擁堵和促進城市可持續發展的關鍵基礎設施. 這一強調高效資源利用與動態供需匹配的交通服務模式，不僅是城市居民便捷出行的日常工具，更是城市交通規劃智慧、運營管理效率與系統韌性的綜合體現. 隨著使用者規模持續增長與運營區域不斷擴張，共用單車系統的供需平衡性、服務公平性、運營經濟性及可持續性，已日益成為運營商和城市管理者必須解決的核心挑戰. 

根據公開的行業資料與研究報導，共用單車日均使用量在過去五年間呈現顯著增長. 與此同時，用戶對車輛可用性（尤其在高峰時段和熱點區域）、調度及時性以及車輛分佈合理性的回饋也同步增多. 這些現象凸顯了對共用單車投放策略與動態調度進行精細化、資料驅動建模的迫切需求. 

在此背景下，一套完善的車輛投放與調度優化系統必須被開發，以確保其在目標城市的服務網路能夠在車輛供給的公平性（覆蓋不同區域需求）、運營的可持續性（降低空駛調度成本）和用戶體驗（提高找車和還車便利性）等方面達到預期目標. 通過科學的投放規劃和高效的調度策略，以最大程度滿足不同區域、時段的出行需求，優化資源利用效率，並提升整體系統的服務水準和經濟效益. 

由於城市的不同區域（如核心商務區、住宅區、交通樞紐、大學城等）具有差異化的出行需求和時空分佈特徵，考慮到系統的運營效率、使用者滿意度和成本控制，我們希望找到一個能夠在城市整體框架內盡可能滿足多方要求的車輛配置與調度方案. 這需要在權衡初始投放成本、日常調度成本、車輛周轉率、區域覆蓋均衡性以及高峰時段回應速度等因素之後，建立一個能夠適應動態需求變化的數學模型. 該模型既要確保各區域使用者享有相對公平的用車機會，也要保證運營商能夠實現可持續的盈利模式和資源的高效利用. 

下文討論了我們需要面對的核心問題：
\begin{enumerate}
    \item 我們如何為城市的不同功能區（或網格區域）確定初始的共用單車投放量？
    在綜合考慮不同區域的關鍵因素（如人口密度、通勤流特徵、POI分佈、公共交通接駁需求、歷史騎行資料、停放設施容量等）之後，我們需要為整個城市劃定服務網格（或功能區），並為每個單元確定一個既能滿足基礎需求、又避免過度投放造成淤積的初始車輛配置方案. 
    \item 我們如何建立一個用於日常（特別是高峰時段）車輛動態調度優化的模型？
    我們的模型需要在現實運營約束（如調度車輛容量、道路通行條件、調度時間視窗、人力成本）下可行，並趨向于高效和低成本. 這需要優化調度路線（最小化總行駛距離/時間）、調度頻次、調度車輛數量，並精准預測各區域未來短時（如未來1至2小時）的車輛供需缺口（短缺或淤積），以動態平衡各區域車輛供給. 
    \item 我們如何使建立的投放與調度模型具有通用性和可擴展性？
    經過通用化設計的模型應當能適應不同規模的城市、不同類型的共用單車（如普通單車、電助力單車）以及不同的運營策略（如是否設置電子圍欄、是否分區計價）. 模型的關鍵參數（如需求預測因數、調度成本係數、公平性權重）應易於根據具體場景進行調整. 
\end{enumerate}

